\chapter*{Реферат}

Итоговая аттестационная работа состоит из 154 страниц, 3 рисунков.

РЕКОМЕНДАТЕЛЬНАЯ СИСТЕМА ДЛЯ ОНЛАЙН-КИНОТЕАТРА НА ОСНОВЕ МЕТОДОВ МАШИННОГО ОБУЧЕНИЯ.

Итоговая аттестационная работа выполнена в формате IT-проекта на тему «Разработка рекомендательной системы для онлайн-кинотеатра с использованием методов машинного обучения». Проект направлен на создание платформы, которая объединяет возможности хранения, обработки данных и персонализации рекомендаций для пользователей. Основная цель проекта — упростить процесс поиска фильмов и предоставить пользователям персонализированные рекомендации на основе их предпочтений.

Дипломный проект основан на использовании современных технологий машинного обучения, таких как коллаборативная фильтрация и контентная фильтрация, которые позволяют точно и быстро анализировать большие объемы данных о фильмах и пользователях. В результате была разработана система, способная генерировать персонализированные рекомендации для пользователей, основываясь на их оценках и предпочтениях, а также на содержании фильмов.

Для достижения поставленной цели были проведены следующие действия: сбор и анализ данных, разработка и обучение моделей машинного обучения, их оценка и улучшение, интеграция моделей в бэкенд-систему, написание тестов и анализ результатов. Основное содержание работы состояло в разработке двух моделей машинного обучения: UserKNN для персонализированных рекомендаций на основе пользовательских оценок и Content-based для рекомендаций на основе содержимого фильмов.

Основным результатом работы, полученным в процессе разработки, является рекомендательная система, оптимизированная для работы в режиме реального времени, написанная на Python с использованием библиотек Django, Django REST Framework, Scikit-learn, Pandas, NumPy и других. Система поддерживает CRUD-операции для управления данными о фильмах, жанрах, режиссёрах и других сущностях, а также предоставляет API для взаимодействия с фронтендом и модулем машинного обучения.

Данные результаты разработки предназначены для использования в онлайн-кинотеатрах и потоковых платформах, которые стремятся улучшить пользовательский опыт за счет персонализированных рекомендаций. Разработку можно адаптировать под нужды различных компаний, занимающихся распространением контента, с помощью интеграции модели с существующими системами.

Применение этого решения позволяет организациям повысить удовлетворённость пользователей, увеличить вовлечённость и улучшить качество рекомендаций. Система может быть модифицирована для работы с различными типами контента, такими как сериалы, музыка или книги, что делает её универсальным инструментом для рекомендательных систем.

Основные результаты:
\begin{enumerate}
    \item Разработана и внедрена рекомендательная система на основе моделей UserKNN и Content-based;
    \item Реализован бэкенд с поддержкой CRUD-операций и API для взаимодействия с фронтендом;
    \item Проведено тестирование системы, подтвердившее её работоспособность и высокую точность рекомендаций
\end{enumerate}

Проект демонстрирует потенциал использования современных технологий машинного обучения для создания эффективных рекомендательных систем, которые могут быть применены в различных отраслях, связанных с обработкой и анализом больших объемов данных.

\endinput
